\documentclass[11pt,letterpaper]{article}
\usepackage[margin=1in]{geometry}
\usepackage{enumitem}
\usepackage{xcolor}
\usepackage{listings}
\usepackage{tcolorbox}
\usepackage{parskip}
\usepackage{titlesec}
\usepackage{fancyvrb}

\definecolor{promptblue}{RGB}{50,100,150}
\definecolor{examplegray}{RGB}{245,245,245}
\definecolor{warnred}{RGB}{180,60,60}
\definecolor{successgreen}{RGB}{60,130,60}

\titleformat{\section}{\large\bfseries\color{promptblue}}{\thesection}{1em}{}
\titleformat{\subsection}{\normalsize\bfseries}{\thesubsection}{1em}{}

\tcbset{
    promptbox/.style={
        colback=examplegray,
        colframe=promptblue,
        boxrule=0.5pt,
        arc=2pt,
        left=6pt,
        right=6pt,
        top=6pt,
        bottom=6pt,
        fontupper=\ttfamily\small
    },
    warnbox/.style={
        colback=red!5,
        colframe=warnred,
        boxrule=0.5pt,
        arc=2pt,
        left=6pt,
        right=6pt,
        top=6pt,
        bottom=6pt
    },
    successbox/.style={
        colback=green!5,
        colframe=successgreen,
        boxrule=0.5pt,
        arc=2pt,
        left=6pt,
        right=6pt,
        top=6pt,
        bottom=6pt
    }
}

\begin{document}

\begin{center}
{\LARGE\textbf{Forcing LaTeX Output from LLMs}}\\[0.5em]
{\large A Prompt Engineering Guide}\\[0.3em]
{\small Preserving creative intent while specifying technical format}
\end{center}

\vspace{1em}
\hrule
\vspace{1em}

\section{The Core Problem}

You want two things simultaneously:

\begin{enumerate}
    \item \textbf{Creative/contextual output} --- the LLM should think, invent, maintain voice, preserve narrative intent
    \item \textbf{Strict formatting} --- output must be valid, compilable LaTeX
\end{enumerate}

These goals can conflict. Over-specifying format can flatten creativity. Under-specifying format gets you markdown or plaintext with \verb|\textbf{}| sprinkled in randomly.

\section{The Key Insight}

\textbf{Separate the WHAT from the HOW.}

Structure your prompt in layers:

\begin{enumerate}
    \item Creative/contextual instruction (what to generate)
    \item Format specification (how to encode it)
    \item Output wrapper (force the container)
\end{enumerate}

\section{Prompt Templates}

\subsection{Template 1: The Wrapper Method}

Wrap your creative prompt, then specify output format as a \textit{container} constraint.

\begin{tcolorbox}[promptbox]
[Your full creative prompt here --- character, setting, \\
intent, style, all the context you've built up]\\
\\
-{}-{}-\\
\\
Generate this in LaTeX. Output ONLY the complete .tex file, \\
starting with \textbackslash documentclass and ending with \\
\textbackslash end\{document\}. No explanation, no markdown \\
code fences, no preamble. Just the raw LaTeX.
\end{tcolorbox}

\begin{tcolorbox}[successbox]
\textbf{Why it works:} The creative instruction gets full attention first. The format spec comes as a \textit{delivery mechanism}, not a constraint on the content itself.
\end{tcolorbox}

\subsection{Template 2: The Role + Format Method}

Assign a role that \textit{implies} LaTeX fluency, then make format feel natural.

\begin{tcolorbox}[promptbox]
You are a writer who drafts everything in LaTeX because you \\
love typographic control. You think in LaTeX---section breaks, \\
emphasis, structured lists are how your mind organizes ideas.\\
\\
[Creative prompt here]\\
\\
Draft this as you naturally would---as a .tex file ready to compile.
\end{tcolorbox}

\begin{tcolorbox}[successbox]
\textbf{Why it works:} The format becomes \textit{characterization} rather than constraint. The LLM isn't ``translating to LaTeX''---it's ``being someone who writes in LaTeX.''
\end{tcolorbox}

\subsection{Template 3: The Specific Structure Method}

When you want particular LaTeX features (colors, boxes, specific packages), specify them as \textit{tools available} rather than requirements.

\begin{tcolorbox}[promptbox]
[Creative prompt]\\
\\
Output as LaTeX with the following available to you:\\
- xcolor for emphasis (define your own semantic colors)\\
- tcolorbox for callouts or asides\\
- enumitem for flexible lists\\
- titlesec if you want styled headers\\
\\
Use whatever serves the content. Full .tex file, documentclass \\
through end\{document\}.
\end{tcolorbox}

\begin{tcolorbox}[successbox]
\textbf{Why it works:} Framing packages as ``available tools'' rather than ``required elements'' gives the LLM creative latitude while ensuring it knows what's possible.
\end{tcolorbox}

\section{Anti-Patterns (What NOT to Do)}

\begin{tcolorbox}[warnbox]
\textbf{Don't: Over-specify structure before content}

\begin{Verbatim}[fontsize=\small]
Write a LaTeX document with:
- 3 sections
- 2 subsections per section  
- A tcolorbox every 200 words
- [creative stuff here]
\end{Verbatim}

This makes the LLM think about \textit{structure} first, \textit{content} second. You get compliant but lifeless output.
\end{tcolorbox}

\begin{tcolorbox}[warnbox]
\textbf{Don't: Mix markdown and LaTeX instructions}

\begin{Verbatim}[fontsize=\small]
Write this as **LaTeX** with `\section{}` headers...
\end{Verbatim}

Mixing markup languages in the prompt confuses the output parser. The LLM may produce hybrid garbage.
\end{tcolorbox}

\begin{tcolorbox}[warnbox]
\textbf{Don't: Forget the document wrapper}

If you just say ``output in LaTeX,'' you might get fragments---a bunch of \verb|\section{}| and \verb|\textbf{}| without the documentclass, packages, or begin/end document. Always specify ``complete .tex file'' or ``compilable document.''
\end{tcolorbox}

\section{Preserving Context Across Format Shifts}

When you've built up context (like our Persistence conversation), and you want to shift to LaTeX output mid-conversation:

\subsection{The Pivot Prompt}

\begin{tcolorbox}[promptbox]
<angle brackets for meta-commentary, not part of output>\\
\\
Take everything we've developed---the world, the tone, the \\
themes, the unresolved questions---and generate [specific thing].\\
\\
Output as LaTeX. The format should SERVE the content: use \\
structure to organize thinking, emphasis to mark importance, \\
color/boxes if they help distinguish idea-types.\\
\\
Complete .tex file. No explanation outside the document.
\end{tcolorbox}

\subsection{Why Angle Brackets / Meta-Commentary}

Using a consistent marker for ``this is instruction, not content'' helps the LLM distinguish:

\begin{itemize}
    \item \texttt{<like this>} --- meta-level, about the output
    \item Everything else --- content-level, part of the output
\end{itemize}

Your original prompt did this intuitively:

\begin{tcolorbox}[promptbox]
< beautiful, thanks>\\
<Can you just make a brainstorm style crude note dump...>
\end{tcolorbox}

The angle brackets signaled: ``this is me talking \textit{about} what I want, not dialogue to include.''

\section{Advanced: Encoding Creative Constraints in LaTeX}

You can use LaTeX structure itself to \textit{embody} creative constraints:

\begin{tcolorbox}[promptbox]
Generate a document where:\\
\\
- \textbackslash section\{\} = major thematic shifts\\
- \textbackslash subsection\{\} = sub-arguments or variations\\
- \textcolor{blue}{blue text} = unresolved questions\\
- \textcolor{red}{red text} = warnings/dangers/things to avoid\\
- \textcolor{green}{green text} = promising directions\\
- tcolorbox = ideas that need isolation/emphasis\\
\\
The structure IS the thinking. Let the LaTeX organization \\
reflect how the ideas relate.
\end{tcolorbox}

This turns LaTeX from ``output format'' into ``thinking tool''---the LLM uses the structure to \textit{organize} its generation, not just \textit{encode} it.

\section{Quick Reference: The Minimal Effective Prompt}

If you just need LaTeX output and don't want to overthink it:

\begin{tcolorbox}[promptbox]
[Your creative prompt, complete with all context]\\
\\
Output as a complete, compilable LaTeX document. \\
Start with \textbackslash documentclass, end with \textbackslash end\{document\}. \\
No markdown, no explanation, no code fences. Just the .tex file.
\end{tcolorbox}

That's usually enough. The key elements:

\begin{enumerate}
    \item ``Complete, compilable'' --- not fragments
    \item ``Start with... end with...'' --- explicit boundaries
    \item ``No markdown, no explanation'' --- prevent contamination
    \item ``Just the .tex file'' --- reinforces raw output
\end{enumerate}

\section{Why This Worked in Our Conversation}

Your prompt:

\begin{tcolorbox}[promptbox]
<Can you just make a brainstorm style crude note dump of \\
ideas and style choices and caveats and directions, you feel \\
me. Generate in LaTeX>
\end{tcolorbox}

Effective because:

\begin{enumerate}
    \item \textbf{Context was pre-loaded} --- we'd built the world together
    \item \textbf{``Brainstorm style crude note dump''} --- gave permission for messy/authentic thinking
    \item \textbf{``you feel me''} --- signaled trust in my interpretation
    \item \textbf{``Generate in LaTeX''} --- clear format spec, came AFTER the creative instruction
\end{enumerate}

The creative latitude came first. The format was the delivery vehicle.

\vspace{2em}
\hrule
\vspace{1em}

\begin{center}
\textit{The format should serve the content.}\\
\textit{Specify what you want made, then how you want it wrapped.}\\
\textit{Trust the LLM to use the tools if you make them available.}
\end{center}

\end{document}
